\scanpage{62}\mspage{61}
mais je ne veux pas l'expliciter en
axiome à part ici, car je le vois comme
cas particulier d'un \uline{axiome de recollement \unsure{des|de} raffinements}, beaucoup
plus général, dont il sera question
plus bas. C'est un axiome qui s'énonce
sans référence particulière aux lieux.
Mais j'ai envie de \unsure{hasarder} deux
axiomes, impliquant encore les lieux.

% slanted
\margin{\unsure{Critère|Critére} de disjonction par les lieux.}
\uline{At} 9 Soient $X$ et $Y\in M$, telles que
$\forall x\in\operatorname{\underline{o}mb}(X),\ y\in\operatorname{\unsure{Omb|omb}}(Y)$, $x$ et $y$ soient
disjoints : \deleted{\unsure{i.e.} $x \mathrel{\tikz[baseline=-.45ex,x=1em,y=1em]{\draw[line width=.055em] (.52,.27)--(.05,.055)--(.48,-.11); \draw[line width=.055em] (.14,.21)--(.46,.04)--(0,-.17);}} y$.} $x\parallel y$. Alors \inserted{et alors \unsure{seult|seulement}} $X$ et $Y$
le sont $\quad X\parallel Y$.

D'autre part, pour un atelier ensembliste,
(et pas seulement \unsure{quasiensembliste|quasi ensembliste}!) on pose

\[
\text{\uline{Atens 2}\quad Si }\ x,y\in L,\ x\ne y,\quad \text{alors }\ x\underset{M}{\mathrel{\tikz[baseline=-.45ex,x=1em,y=1em]{\draw[line width=.055em] (.52,.27)--(.05,.055)--(.48,-.11); \draw[line width=.055em] (.14,.21)--(.46,.04)--(0,-.17);}}}y
\]
qui se traduit en termes de \deleted{$\unsure{M|L}, \mathrel{\tikz[baseline=-.45ex,x=1em,y=1em]{\draw[line width=.055em] (.52,.27)--(.05,.055)--(.48,-.11); \draw[line width=.055em] (.14,.21)--(.46,.04)--(0,-.17);}}$}
$\unsure{.|-}\,M\subset\operatorname{Figel}(\deleted{\unsure{L|l}}L)$, et $\underset{M}{\mathrel{\tikz[baseline=-.45ex,x=1em,y=1em]{\draw[line width=.055em] (.52,.27)--(.05,.055)--(.48,-.11); \draw[line width=.055em] (.14,.21)--(.46,.04)--(0,-.17);}}}$, par
\[
\text{\uline{Atens\deleted{\unsure{é|*}}\,\deleted{\unsure{g|s}}\unsure{tf|t.f|fig} 4}}\qquad
\forall\,x\deleted{\unsure{\in|\,\in L}},y\in L,\ x\ne y,\quad \text{on a }\{\{x\}\}\underset{M}{\mathrel{\tikz[baseline=-.45ex,x=1em,y=1em]{\draw[line width=.055em] (.52,.27)--(.05,.055)--(.48,-.11); \draw[line width=.055em] (.14,.21)--(.46,.04)--(0,-.17);}}}\{\{y\}\}.
\]
